\chapter{Conclusiones y trabajo futuro}
\label{chap:conclusiones}

\ph{El último capítulo cierra la memoria en tres pasos: las
\emph{conclusiones} (qué se ha conseguido, apoyado en la evidencia del
Capítulo~\ref{chap:resultados}, que el lector acaba de leer), las
\emph{limitaciones} (qué no se ha conseguido o no se ha podido validar) y
el \emph{trabajo futuro} (qué líneas abre todo lo anterior).}

\section{Conclusiones}
\label{sec:conclusiones}

\ph{Estructura recomendada:}

\ph{1. \textbf{Recapitulación del problema}: una o dos frases recordando
qué problema se planteaba al inicio (sin entrar en detalles, ya que el
lector viene del Capítulo~\ref{chap:resultados}).}

\ph{2. \textbf{Cumplimiento de los objetivos}: revisa los objetivos
específicos planteados en el Capítulo~\ref{chap:intro} y argumenta cómo
se han cubierto, \textbf{apoyando cada afirmación en evidencia concreta}:
la tabla, figura o sección del Capítulo~\ref{chap:resultados} que lo
demuestra (o, si procede, de los Capítulos~\ref{chap:requisitos}
a~\ref{chap:pruebas}). Evita el «se ha cumplido» sin soporte: todo
cumplimiento debe poder verificarse en un capítulo anterior. Dedica un
párrafo a cada objetivo, en el mismo orden en que se enunciaron.}

\begin{itemize}
  \item \ph{El \textbf{primer objetivo} \dots se ha cumplido mediante
        \dots, como muestra la Tabla/Figura~\dots\ del
        Capítulo~\ref{chap:resultados}.}
  \item \ph{El \textbf{segundo objetivo} \dots se ha cumplido mediante
        \dots, como muestra \dots}
  \item \ph{El \textbf{tercer objetivo} \dots se ha cumplido mediante
        \dots, como muestra \dots}
  \item \ph{El \textbf{cuarto objetivo} \dots se ha cumplido mediante
        \dots, como muestra \dots}
  \item \ph{El \textbf{quinto objetivo} \dots se ha cumplido
        \emph{parcialmente} mediante \dots, como se discute en la
        Sección~\ref{sec:limitaciones}.}
\end{itemize}

\ph{La siguiente tabla sintetiza el cumplimiento y cierra la cadena
abierta por la Tabla~\ref{tab:criterios-exito} del
Capítulo~\ref{chap:intro}: mismos objetivos, mismos criterios de éxito,
ahora con su grado de cumplimiento y la evidencia que lo acredita.
La columna «Grado» usa la escala Completo / Parcial / No alcanzado:
«Completo» significa que el criterio se cumple íntegramente con la
evidencia citada; «Parcial», que se cumple solo en parte, explicando qué
falta en la Sección~\ref{sec:limitaciones}; «No alcanzado», que el
criterio no se ha cumplido, analizando la causa en la
Sección~\ref{sec:limitaciones}. Admitir un «Parcial» con remisión honesta
a la Sección~\ref{sec:limitaciones} es señal de madurez, no de debilidad.}

\begin{longtable}{@{}P{1cm} P{5.5cm} P{2cm} P{5cm}@{}}
  \toprule
  \textbf{Obj.} & \textbf{Criterio de éxito} & \textbf{Grado} & \textbf{Evidencia (ref.)} \\
  \midrule
  \endhead
  O1 & \ph{Criterio de éxito del objetivo 1.} & \ph{Completo} & \ph{Tabla/Figura/Sección que lo acredita.} \\[0.3em]
  O2 & \ph{Criterio de éxito del objetivo 2.} & \ph{Completo} & \ph{Evidencia.} \\[0.3em]
  O3 & \ph{Criterio de éxito del objetivo 3.} & \ph{Completo} & \ph{Evidencia.} \\[0.3em]
  O4 & \ph{Criterio de éxito del objetivo 4.} & \ph{Completo} & \ph{Evidencia.} \\[0.3em]
  O5 & \ph{Criterio de éxito del objetivo 5.} & \ph{Parcial}  & \ph{Evidencia y remisión a la Sección~\ref{sec:limitaciones}.} \\
  \bottomrule
  \caption{Grado de cumplimiento de los objetivos y evidencia asociada.}
  \label{tab:cumplimiento-objetivos}
\end{longtable}

\ph{3. \textbf{Aportaciones más amplias}: comenta brevemente la contribución
del trabajo más allá del producto en sí. Por ejemplo, aplicación de un
paradigma a un dominio nuevo, integración con un ecosistema, validación
empírica de una técnica, etc.}

\ph{4. \textbf{Reflexión crítica y lecciones aprendidas}: resume en un
párrafo la desviación global de horas y en qué fases se concentró,
remitiendo al cierre del Capítulo~\ref{chap:planificacion}
(Sección~\ref{sec:cierre}) para el detalle, sin duplicarlo. Añade aquí
una o dos lecciones \emph{nuevas}, técnicas o personales: qué decisión
repetirías, qué harías de otra forma si empezaras hoy, qué sabes ahora
que no sabías al comenzar. Dos o tres frases honestas valen más que un
párrafo genérico.}

\ph{5. \textbf{Cierre}: una frase que cierre el capítulo conectando las
limitaciones identificadas en la introducción con la solución desarrollada
(``el sistema desarrollado resuelve los X problemas estructurales\dots'').}

El objetivo principal de este trabajo era desarrollar un generador automático de
modelos de características capaz de producir modelos UVL con diferentes niveles
de complejidad y expresividad, integrándose dentro del ecosistema Flamapy y
UVLHub. Para ello, se ha desarrollado una solución configurable que permite
generar modelos con construcciones avanzadas del lenguaje UVL y validar su
correcto funcionamiento mediante pruebas automatizadas.

Los resultados obtenidos en el Capítulo~\ref{chap:resultados} permiten concluir
que los objetivos planteados inicialmente han sido alcanzados. El objetivo O1,
relacionado con la generación automática de modelos UVL con distintos niveles de
expresividad, se ha cumplido mediante la generación y validación de modelos con
diferentes configuraciones, como se muestra en la Sección~\ref{subsec:resultados_o1}.

El objetivo O2 se ha cumplido mediante la incorporación de un sistema de
configuración flexible que permite adaptar el proceso de generación a distintos
escenarios. Por su parte, el objetivo O3 se ha alcanzado mediante la integración
del generador como plugin de Flamapy y su utilización desde UVLHub, permitiendo
su incorporación dentro del flujo habitual de la plataforma
(Secciones~\ref{subsec:resultados_o2} y~\ref{subsec:resultados_o3}).

El objetivo O4, centrado en la evaluación del rendimiento del generador, se ha
verificado mediante la medición del coste de generación en escenarios con un
volumen elevado de modelos, obteniendo resultados adecuados para el uso
previsto de la herramienta (Sección~\ref{subsec:resultados_o4}). Finalmente, el
objetivo O5 se ha cumplido mediante la ejecución de una batería de pruebas
automatizadas que valida las principales funcionalidades del generador y sus
diferentes flujos de configuración
(Sección~\ref{subsec:resultados_o5}).

La Tabla~\ref{tab:grado-cumplimiento-objetivos} resume el grado de cumplimiento
de los objetivos definidos.

\begin{table}[H]
\centering
\small
\begin{tabular}{@{}P{1cm} P{6cm} P{2cm}@{}}
\toprule
\textbf{Obj.} & \textbf{Objetivo} & \textbf{Grado} \\
\midrule
O1 & Generación automática de modelos UVL con expresividad avanzada. & Completo \\
O2 & Configuración flexible del proceso de generación. & Completo \\
O3 & Integración con Flamapy y UVLHub. & Completo \\
O4 & Evaluación temporal del generador. & Completo \\
O5 & Validación mediante pruebas automatizadas. & Completo \\
\bottomrule
\end{tabular}
\caption{Grado de cumplimiento de los objetivos. Elaboración propia}
\label{tab:grado-cumplimiento-objetivos}
\end{table}

Como aportación principal, este trabajo proporciona una solución integrada para
la generación automática de modelos UVL con mayor expresividad que los
generadores tradicionales centrados en construcciones booleanas. Además, la
integración dentro del ecosistema Flamapy y UVLHub permite utilizar esta
funcionalidad en escenarios de experimentación y análisis de variabilidad.

A nivel personal y técnico, el desarrollo ha permitido identificar la
importancia de realizar una planificación temporal con margen para imprevistos,
así como la necesidad de utilizar desde el inicio herramientas que faciliten un
entorno reproducible, como Docker. También se ha comprobado la relevancia de
realizar refactorizaciones progresivas y acompañar los cambios funcionales con
pruebas automatizadas para evitar la acumulación de errores durante la
evolución del software.

En conjunto, los resultados obtenidos permiten concluir que la solución
desarrollada cumple los objetivos establecidos y proporciona una base extensible
para futuras mejoras relacionadas con la generación y análisis de modelos de
características.

\section{Limitaciones}
\label{sec:limitaciones}

\ph{Reflexión honesta sobre lo que el trabajo \emph{no} hace o no ha
podido demostrar. Reconocer limitaciones no resta nota: demuestra que
comprendes el alcance real de lo construido. Para cada limitación sigue
el patrón \emph{causa} (decisión, restricción de tiempo o de recursos)
+ \emph{efecto} práctico + \emph{por qué es asumible o cómo se
mitigaría}. Organízalas en tres categorías:}

\begin{itemize}
  \item \textbf{\ph{Limitación técnica:}} \ph{qué no hace el sistema o
        dónde rinde peor: p.\,ej., rendimiento no evaluado bajo carga
        real, dependencia de un servicio externo, escalabilidad no
        contrastada\dots\ Causa, efecto y mitigación.}
  \item \textbf{\ph{Limitación de validación:}} \ph{qué no se ha podido
        medir o con qué amenazas a la validez: p.\,ej., datos sintéticos,
        sin usuarios finales, una sola configuración de prueba\dots\
        Conecta con la discusión de la Sección~\ref{sec:discusion} del
        Capítulo~\ref{chap:resultados}.}
  \item \textbf{\ph{Limitación de alcance:}} \ph{qué quedó fuera del
        trabajo y por qué, de forma coherente con lo declarado en la
        Sección~\ref{sec:alcance}.}
\end{itemize}

\ph{Cada limitación es candidata natural a una línea de la
Sección~\ref{sec:trabajo_futuro}: si una limitación no genera línea
futura ni se justifica como asumible, queda coja.}

A pesar de los resultados obtenidos, existen ciertas limitaciones que deben
tenerse en cuenta al interpretar el alcance del trabajo realizado. Estas
limitaciones no afectan al cumplimiento de los objetivos planteados, pero
delimitan los escenarios en los que se han validado las capacidades del
generador.

Desde un punto de vista técnico, la evaluación del generador se ha realizado
sobre configuraciones representativas, pero no se ha estudiado su comportamiento
en escenarios extremos con modelos de características de gran tamaño o con una
complejidad significativamente superior. Esto puede afectar al coste
computacional en configuraciones con un número elevado de características,
atributos o restricciones. No obstante, la arquitectura desarrollada permite
continuar optimizando y ampliando el generador para este tipo de escenarios.

Respecto a la validación, las pruebas realizadas se han basado principalmente
en modelos sintéticos generados mediante la propia herramienta y en escenarios
diseñados para cubrir las funcionalidades implementadas. Por tanto, aunque la
suite de pruebas automatizadas permite validar el comportamiento esperado del
sistema, no se ha evaluado todavía su comportamiento con modelos reales
procedentes de dominios específicos ni mediante estudios con usuarios finales.
Además, no se ha realizado una comparación cuantitativa directa con otras
herramientas debido a la ausencia de condiciones experimentales homogéneas que
permitieran obtener resultados comparables.

Finalmente, en cuanto al alcance del trabajo, la solución desarrollada se ha
centrado en la generación automática de modelos UVL y su integración dentro del
ecosistema Flamapy y UVLHub. No se han abordado otras líneas relacionadas, como
la optimización avanzada del proceso de generación, nuevas estrategias de
búsqueda o análisis posterior de los modelos generados, ya que quedan fuera de
los objetivos definidos inicialmente.

\section{Trabajo futuro}
\label{sec:trabajo_futuro}

\ph{Líneas de extensión del trabajo, ordenadas de mayor a menor
inmediatez. Cada línea debe responder a tres preguntas: \emph{qué} se
haría, \emph{por qué} merece la pena (anclándola a una limitación de la
Sección~\ref{sec:limitaciones} o a una tendencia del
Capítulo~\ref{chap:estado_arte}, idealmente con cita) y qué
\emph{viabilidad o dificultad} estimada tiene. Al menos una línea debe
ser genuinamente innovadora (no «añadir más tests»), explicando qué la
hace novedosa.}

\textbf{\ph{Línea 1 (corto plazo): título corto.}} \ph{Qué: \dots\
Por qué: resuelve la limitación técnica identificada en la
Sección~\ref{sec:limitaciones}. Viabilidad: encaja con la arquitectura
actual porque \dots; dificultad estimada baja.}

\textbf{\ph{Línea 2 (corto plazo): título corto.}} \ph{Qué / por qué
(limitación de validación) / viabilidad-dificultad.}

\textbf{\ph{Línea 3 (medio plazo): título corto.}} \ph{Qué / por qué
(funcionalidad excluida del alcance) / viabilidad-dificultad.}

\textbf{\ph{Línea 4 (innovadora): título corto.}} \ph{Qué: extensión
ambiciosa (nueva técnica, dominio o integración). Por qué: conecta con
la tendencia identificada en el estado del
arte~\cite{apellido2024articulo}. Qué la hace novedosa: \dots\
Viabilidad: qué haría falta para abordarla y por qué este trabajo es un
buen punto de partida.}

A partir de las limitaciones identificadas y de las posibilidades de extensión
del sistema desarrollado, se plantean diferentes líneas de trabajo futuro
orientadas a mejorar la reutilización de configuraciones, la integración con
UVLHub y la capacidad del generador para adaptarse a nuevos escenarios de uso.

\subsubsection*{Exportación e importación de configuraciones}

Una primera línea de trabajo consistiría en desarrollar una transformación de
Flamapy que permita exportar la configuración seleccionada del generador en un
formato estándar como JSON, permitiendo posteriormente importar dicha
configuración desde UVLHub para reproducir una generación concreta.

Esta extensión mejoraría la reproducibilidad de los experimentos y facilitaría
el intercambio de configuraciones entre usuarios, evitando la necesidad de
definir manualmente conjuntos complejos de parámetros. Además, permitiría
compartir escenarios de generación de forma sencilla dentro del ecosistema de
herramientas de variabilidad.

La implementación de esta funcionalidad presenta una viabilidad elevada, ya que
los parámetros del generador ya se encuentran estructurados internamente. Sería
necesario únicamente definir un esquema estable de serialización y añadir los
mecanismos de exportación e importación correspondientes.

\subsubsection*{Publicación directa de modelos en UVLHub}

Otra posible extensión consistiría en ampliar la integración con UVLHub para
permitir subir directamente los modelos generados a la plataforma, pudiendo
establecer su visibilidad y compartirlos con otros usuarios.

Actualmente, el generador permite producir modelos UVL, pero la gestión posterior
de estos modelos requiere pasos adicionales. La incorporación de esta
funcionalidad permitiría centralizar el flujo completo, desde la generación
hasta el almacenamiento y distribución de modelos de características.

Esta línea presenta una viabilidad media-alta, ya que requeriría ampliar la
comunicación entre el generador y los servicios existentes de UVLHub, además de
definir aspectos relacionados con permisos, almacenamiento y gestión de modelos
públicos.

\subsubsection*{Evaluación con modelos reales y benchmarking}

Una tercera línea de trabajo consistiría en realizar una evaluación experimental
del generador utilizando modelos de características procedentes de dominios
reales y comparando sus resultados con otras herramientas existentes bajo unas
condiciones experimentales comunes.

Esta extensión permitiría superar una de las limitaciones identificadas en la
evaluación actual, basada principalmente en modelos sintéticos, y proporcionaría
una medida más precisa de la utilidad del generador en escenarios reales.

Su viabilidad es media, ya que requeriría disponer de conjuntos de datos
adecuados, establecer una metodología experimental común y adaptar las
herramientas comparadas para garantizar unas condiciones equivalentes.

\subsubsection*{Optimización del proceso de generación}

Una posible línea de trabajo futuro consistiría en optimizar el proceso interno
de generación para mejorar el rendimiento en configuraciones con una mayor
complejidad estructural o un volumen elevado de modelos generados.

Aunque los resultados obtenidos muestran un comportamiento adecuado en los
escenarios evaluados, el coste computacional puede incrementarse al aumentar el
número de características, atributos o restricciones. Por ello, sería
interesante analizar los puntos críticos del proceso de generación y aplicar
técnicas de optimización, como mejoras en la gestión de estructuras internas,
reducción de operaciones redundantes o paralelización de determinadas etapas.

Esta extensión presenta una viabilidad media-alta, ya que no requiere modificar
la arquitectura general del sistema, sino realizar mejoras sobre los componentes
existentes del motor de generación. La principal dificultad reside en identificar
aquellas operaciones que tienen mayor impacto sobre el tiempo de ejecución y
mantener la corrección de los modelos generados tras las optimizaciones.

\subsubsection*{Generación guiada mediante aprendizaje automático}

Como línea de investigación más ambiciosa, se plantea la incorporación de
técnicas de aprendizaje automático para aprender patrones presentes en modelos
de características reales y utilizarlos durante el proceso de generación.

Esta aproximación permitiría evolucionar desde una generación basada únicamente
en parámetros definidos manualmente hacia una generación capaz de producir
modelos sintéticos adaptados a las características de dominios concretos. De
esta forma, se podrían obtener modelos más realistas para tareas de evaluación
y experimentación.

La viabilidad de esta línea es más reducida, ya que requeriría disponer de un
volumen suficiente de modelos reales, definir representaciones adecuadas de los
modelos de características y desarrollar técnicas capaces de capturar sus
patrones estructurales. No obstante, constituye una posible evolución innovadora
del generador desarrollado.